\subsection{标准正交化}

定理 2. --- 对希尔伯特空间 E 中的每个标准正交集 L，存在 E 的包含 L 的标准正交基 B。

事实上，设 D 为 E 的所有标准正交子集按包含关系线性排序所成的族；立即可以看出此族具有有限特征（S, III, § 4, No.~5）。于是由 S, III, § 4, No.~5 的定理 1，D 中存在包含 L 的极大族 B。剩下要证 B 是完全集。若不然，将存在一个与 B 的所有向量都正交的向量 \(y \neq 0\)（V, p.~22, 命题 5），用适当的标量乘 y，可设 \(\|y\| = 1\)；那么 \(B \cup \{y\}\) 将是异于 B 且包含 B 的标准正交集；这与 B 的定义矛盾；定理得证。

推论 1. --- 在每个希尔伯特空间中都存在标准正交基。

只需对 L = ∅ 的情形应用定理 2。

推论 2. --- 每个希尔伯特空间都同构于某个空间 \(\ell^2 (\mathbf{I})\)。

更确切地说，设 \((e_{i})_{i\in I}\) 为希尔伯特空间 E 的标准正交基。由命题 4（V, p.~21）与命题 5（V, p.~22），由下式定义的映射 \(\phi\)

\[
\phi (x) = \big (\langle e _ {i} | x \rangle \big) _ {i \in \mathrm{I}}\tag{4}
\]

是从 E 到 \(\ell^2 (\mathrm{I})\) 上的希尔伯特空间同构。逆同构 \(\psi\) 由下式定义

\[
\psi ((\lambda_ {i}) _ {i \in \mathbf {I}}) = \sum_ {i \in \mathbf {I}} \lambda_ {i} e _ {i}.\tag{5}
\]

命题 6. --- 设 E 为分离的准希尔伯特空间，\((a_{n})_{n\in\mathbb{I}}\)（I 是 N 的以 1 为起点的区间）为 E 中向量的一个可数的（有限或无限）无关族。则存在唯一的 E 中的标准正交族 \((e_{n})_{n\in\mathbb{I}}\)，具有下列性质：

\begin{enumerate}
\def\labelenumi{\arabic{enumi})}
\item
  对每个整数 \(p \in \mathbf{I}\)，\(\mathbf{E}\) 的由 \(e_1, e_2, \ldots, e_p\) 生成的向量子空间与 \(\mathbf{E}\) 的由 \(a_1, a_2, \ldots, a_p\) 生成的向量子空间相同；
\item
  对每个整数 \(p \in \mathbf{I}\)，数 \(\langle a_p | e_p \rangle\) 是实数且 \(> 0\)。
\end{enumerate}

事实上，设 \(V_{n}\) 为由 \(a_{1}, a_{2}, ..., a_{n}\) 生成的（n 维）子空间。若 \(n + 1 \in I\) 且 \(b_{n+1} = a_{n+1} - p_{\mathrm{V}_{n}}(a_{n+1})\)（其中 \(p_{V_{n}}\) 是到完备子空间 \(V_{n}\) 上的正交投影算子），则直线 \(Kb_{n+1}\) 是 \(V_{n}\) 在 \(V_{n+1}\) 中的正交（补）。若诸 \(e_{n}\) 满足命题的条件 1)，则必须有 \(e_{n+1} = \lambda b_{n+1}\)；条件 \(\|e_{n+1}\| = 1\) 于是蕴含 \(|\lambda|^2 \|b_{n+1}\|^2 = 1\)，而条件 \(\langle a_{n+1}|e_{n+1} \rangle > 0\) 蕴含 \(\lambda \langle a_{n+1}|b_{n+1} \rangle > 0\)；这就完全确定了 \(\lambda\)，于是我们证明了可以用归纳法确定唯一的标准正交族 \((e_n)_{n\in I}\)，使之满足命题的条件 1) 与 2)。

称序列 \((e_{n})_{n\in\mathbb{I}}\) 是由无关族 \((a_{n})_{n\in\mathbb{I}}\) 经标准正交化得到的。显然，由族 \((e_{n})\) 生成的向量子空间与由族 \((a_{n})\) 生成的向量子空间相同。特别地，若 \((a_{n})\) 是完全序列，则 \((e_{n})\) 也是，从而它是 E 的标准正交基；由此得到：

推论. --- 在每个满足第一可数公理的分离准希尔伯特空间 E 中，都存在可数的标准正交基。

若 E 满足第一可数公理，则 E 中存在完全序列，而且总能从此序列中抽取一个无关的完全族（A, II, § 7, No.~1, 定理 2）。

可以举出没有任何标准正交基的分离准希尔伯特空间的例子（V, p.~70, 练习 2）。

例. --- 设 I 为 R 的区间 \((-1,1)\)，E 为 I 上实值连续函数的向量空间。用 x 表示从 I 到 R 中的典范单射，把它看作 E 的一个元素。由 Stone-Weierstrass 定理，序列 \((x^{n})_{n\in\mathbf{N}}\) 在 E 中对一致收敛拓扑是完全的（GT, X, § 4, No.~2）。

把 E 看作实的分离准希尔伯特空间，其内积由下式给出

\[
\langle f | g \rangle = \int_ {- 1} ^ {1} f (t) g (t) d t.
\]

于是序列 \((x^{n})_{n\in \mathbb{N}}\) 在准希尔伯特空间 E 中是完全的。设 \((\Pi_n)_{n\in \mathbb{N}}\) 为由序列 \((x^n)_{n\in \mathbb{N}}\) 经标准正交化得到的序列。可以证明 \(\Pi_n = (n + \frac{1}{2})^{1 / 2}\mathbf{P}_n\)，其中 Legendre 多项式 \(\mathbf{P}_n\) 由下式定义

\[
\mathrm{P} _ {n} (x) = \frac {1}{2 ^ {n} n !} \left(\frac {d}{d x}\right) ^ {n} (x ^ {2} - 1) ^ {n}.
\]

命题 7. --- 在希尔伯特空间 E 中，两个标准正交基是等势的。

设 B 与 C 为 E 的两个标准正交基。当两集 B、C 之一为有限时，情形是平凡的，因为有限的标准正交基是空间的代数基。因此设 B 与 C 都是无限的。对每个 \(x \in B\)，设 \(C_{x}\) 为由 \(y \in C\) 中所有使 \(\langle x | y \rangle \neq 0\) 者组成的子集。集 \(C_{x}\) 是可数的（V, p.~21, 命题 4）。对每个 \(y \in C\)，存在 \(x \in B\) 使 \(y \in C_{x}\)，因为 B 是标准正交基且 \(y \neq 0\)；换言之，C 是当 x 取遍 B 时诸可数集 \(C_{x}\) 的并。于是 C 的基数小于 N × B 的基数，从而小于 B 的基数（S, III, §6, No.~3, 推论 4）；类似地，B 的基数小于 C 的基数；证明完毕。

希尔伯特空间 E 的任一标准正交基的基数称为 E 的希尔伯特维数。

推论 1. --- 给定希尔伯特空间 E 中的两个标准正交基，存在 E 的自同构把第一个基变为第二个基。

推论 2. --- 为使希尔伯特空间 \(\ell^2 (\mathbf{I})\) 与 \(\ell^2 (\mathbf{J})\) 同构，必要且充分的是 I 与 J 等势。

